\documentclass[10pt]{article}
\usepackage{template}
\usepackage{amsmath,amssymb,tikz,graphicx,xurl,microtype}
\usetikzlibrary{arrows.meta,positioning,calc}
\setlength{\emergencystretch}{3em}
\raggedbottom
\SetKeys{date=2026-11-02,vol=020,title={每日有趣新闻｜停火的条件与验证的边界},author={ChatGPT 编辑部},note={\small 采编截止：北京时间11月2日上午。主要报道10月31日至11月1日发生或更新的事件；周末arXiv未发布新一期数学列表，数学栏选取10月30日新列表中此前未报道的预印本。所有时政数字均注明发布方与口径。}}
\renewcommand{\newsitem}[3]{\par\medskip{\large\bfseries\color{newsblue}#1\par}{\footnotesize\color{newsgray}#2\par}\smallskip\color{black}#3\par\medskip{\color{newsgray}\rule{\linewidth}{0.25pt}}\par}
\newcommand{\newslead}[3]{\par\medskip{\LARGE\bfseries\color{newsblue}#1\par}\vspace{1mm}{\small\color{newsgray}#2\par}\smallskip\color{black}#3\par\medskip{\color{newsgray}\rule{\linewidth}{0.45pt}}\par}
\newcommand{\newssub}[1]{\par\smallskip\noindent{\bfseries\color{newsblue}#1}\quad}
\begin{document}
\begin{documents}
\raggedcolumns
\newslead{霍尔木兹海峡能在「两天内」重开吗？多哈斡旋进入第三天}{国际政治深度｜11月1日；路透社、AP及伊朗外长表态交叉核查}{
11月1日，伊朗谈判人员在卡塔尔首都多哈连续第三天与调解方接触，试图寻找结束持续八个月的美以对伊战争的协议。路透社引述了解磋商情况的官员说，卡塔尔在与伊朗接触的同时保持与华盛顿的沟通。伊朗外长阿拉格齐确认讨论仍在推进，但尚未达成最终结论。最引人注意的是一项被媒体称为「两日方案」的设想：\textbf{如果美伊达成协议，霍尔木兹海峡可能在随后两天内重新开放}。这是条件句，不是航运已恢复的事实。\cite{ReutersIran,APIran}

\newssub{谈判桌上究竟坐着谁？}目前公开报道能够支持的描述是「伊朗与卡塔尔调解人磋商，卡塔尔与美国沟通」，而不是美伊双方已经恢复面对面的正式谈判。伊朗方面尤其强调不存在直接美伊会谈。卡塔尔既是海湾国家，又是液化天然气出口国，海峡不畅会影响本国及其他能源出口方的商业利益。调解者的利益不必与交战方相同，但也不能据此断言其天然中立。\cite{ReutersIran,APIran}

\begin{center}
\resizebox{\linewidth}{!}{\begin{tikzpicture}[>=Latex,scale=.88,every node/.style={font=\scriptsize}]
\node[draw,rounded corners,fill=blue!7,text width=2.25cm,align=center] (iran) at (0,0) {伊朗谈判人员\\德黑兰／多哈};
\node[draw,rounded corners,fill=orange!10,text width=2.25cm,align=center] (qatar) at (4,0) {卡塔尔调解方\\传递条件与回应};
\node[draw,rounded corners,fill=blue!7,text width=2.25cm,align=center] (usa) at (8,0) {美国决策方\\尚无最终协议};
\draw[<->,thick] (iran) -- node[above]{磋商} (qatar);
\draw[<->,thick] (qatar) -- node[above]{间接沟通} (usa);
\node[draw,dashed,rounded corners,fill=gray!5,text width=6.5cm,align=center] at (4,-1.35) {「协议后两日内重开海峡」是被讨论的\textbf{条件性方案}，不是已经发生的通航恢复};
\end{tikzpicture}}
\par\scriptsize 图1｜编辑部根据11月1日路透社、AP报道绘制的谈判关系示意；箭头表示信息沟通，不代表直接美伊会谈。\end{center}\color{black}

\newssub{为什么是海峡，而不先谈停火？}霍尔木兹海峡是波斯湾通向阿拉伯海的咽喉。它的地理位置让油轮、液化天然气船与普通货船的安全问题直接关联能源价格、海运保险和各国库存。战争造成的损失还包括港口、油气设施、运输链条的停工或损毁；即便航道名义上开放，船东也可能因为保险、船员安全和港口接卸能力而不立即恢复航次。由此，\textbf{「宣布开放」「实际安全通行」「能源供给恢复」是三件不同的事}。\cite{ReutersIran,ReutersOPEC}

阿拉格齐在公开讲话中表示伊朗既准备防卫也准备外交，并将结果系于美国的答复。美国国务卿鲁比奥则称总统尚未决定是否再度对伊朗发动袭击，更倾向外交解决。美国总统特朗普此前承诺11月3日中期选举前不发起新的袭击，但没有排除选举后使用武力。双方表态显示各自仍保留施压空间；它们不等于签署了可执行的停火文本。\cite{ReutersIran,APIran}

\newssub{三类条件不能混成一个口号。}一项可执行的海运安排至少要回答三个问题：其一，通航权的适用范围是所有商船，还是只限获准航线、特定船旗与货类？其二，船只遇袭、被扣押或发生误判时由谁核查并通报？其三，经济制裁、军事部署与地区其他战线之间如何联动？在公开方案全文尚不可核查的情况下，不能把媒体概括的「两日内重开」扩写为伊朗已接受美国全部条件，或美国已承诺撤军、解除制裁。\cite{ReutersIran}

\newssub{现实风险不只来自谈判失败。}美联社11月1日报道，卡塔尔天然气运输公司Nakilat称其一艘船在波斯湾被不明物体击中，船只受损，但无人受伤，也没有报告环境损害；公司没有认定袭击者。卡塔尔外交部谴责事件并表示调查。此处最重要的事实边界是\textbf{尚未确定是谁发射了物体}。把它直接写成某一方的攻击，会把待调查的责任问题伪装成既定事实。\cite{APIran}

\newssub{编辑观察：协议的可验证性比漂亮措辞重要。}停火谈判经常先宣布一个容易传播的期限，却在执行时遭遇定义争议：什么算攻击、谁负责监测、如何核对违约？一份海运协议即使不能结束陆上战事，也可能切实减少商船风险；反过来，一份政治宣言若没有执行机制，也可能只短暂影响市场情绪。接下来值得观察的是卡塔尔或美国是否公布正式条款、伊朗是否公布明确的通航指令、船舶航迹与保险成本是否同步变化，以及任何违约是否能由独立渠道核实。\cite{ReutersIran,ReutersOPEC}
}

\newslead{币安遭窃63亿美元：补上95\%的资产，是否就等于找回了钱？}{网络安全与金融深度｜10月30日至11月1日；路透社、FBI归因及币安声明}{
11月1日，全球大型加密资产交易平台币安分阶段恢复提现。此前，攻击者在约20分钟内从平台相关钱包转走估值约63亿美元的数字资产。路透社称，美国联邦调查局把袭击归因于朝鲜关联的TraderTraitor组织；币安则表示，已经通过公司资金、保险基金与用户保护储备等渠道补足或覆盖约95\%的被盗资产。\textbf{「补足损失」和「追回被盗币」不是同义词}：前者可能是平台自己掏钱买回等值资产，后者才意味着原本被转走的资产回到控制之下。\cite{ReutersBinance,ReutersBinanceCover,ReutersTRM}

\newssub{什么叫「覆盖95\%」？}币安披露已购买或安排约60亿美元替代资产，剩余约5\%主要是流动性不足、短期难以大量买入的币种。公司称已为损失预留约63.4亿美元，并计划恢复其用户安全资产基金（SAFU）的规模。这个数字关系到客户账面资产能否重新得到支持，却并不证明被盗链上资产已全部冻结、攻击者已经被抓获，也不保证未来提现不会遭遇新的拥堵。\cite{ReutersBinance,ReutersBinanceCover}

\begin{center}
\begin{tikzpicture}[x=.9cm,y=.65cm,every node/.style={font=\scriptsize}]
\draw[->] (0,0)--(8,0) node[right]{十亿美元};
\foreach \x in {0,2,4,6}{\draw (\x,0.09)--(\x,-.09) node[below]{\x};}
\fill[blue!45] (0,.65) rectangle (6.3,1.15);
\node[anchor=west] at (6.4,.9) {6.3 被盗资产估值};
\fill[teal!55] (0,1.6) rectangle (6.0,2.1);
\node[anchor=west] at (6.1,1.85) {约6.0 已补足};
\fill[orange!60] (0,2.55) rectangle (.3,3.05);
\node[anchor=west] at (.42,2.8) {约0.3 尚未补足（估算）};
\end{tikzpicture}
\par\scriptsize 图2｜据路透社11月1日报道及币安声明绘制。金额为当时披露的美元估值；「补足」包含购买替代资产，不等于全部被盗币已追回。\end{center}\color{black}

\newssub{交易所为什么比个人钱包更复杂？}区块链上的一笔转账可以公开查询，但「链上公开」不意味着攻击难以得手。大型平台同时管理热钱包、冷钱包、内部账本、跨链资产、客户提现队列及员工访问权限。报道援引币安高管说法称，事件涉及一台被攻陷的笔记本电脑和用于清除证据的恶意软件；这一描述来自企业自身，完整取证报告尚未公开。攻击者获得了哪一类签名权限、内部审批流程何时失效，仍需要独立技术审计，而不能凭一段新闻文字还原入侵链。\cite{ReutersBinance}

\newssub{归因证据也有层次。}FBI声明将行动归于TraderTraitor；TRM Labs等区块链分析公司也提出高度确信的朝鲜关联判断。美国总统特朗普却公开表示尚不确定是否为朝鲜所为。执法机构的归因、商业分析机构的链上追踪、政治人物的表态具有不同证据基础和立场。朝鲜长期否认国家支持此类盗窃。\textbf{目前能准确报道的是「FBI归因于某组织」，而非本刊独立证明朝鲜政府直接下令}。\cite{ReutersBinance,ReutersTRM}

\newssub{为何恢复提现仍是压力测试？}路透社11月1日更新报道，客户在提现恢复后累计提取约36.4亿美元，约占平台报告资产的2.1\%。第三方数据平台DefiLlama给出更高的流出估计，币安则认为其统计没有正确计入转入新热钱包的资金。这里的「提款额」「净流出」「内部钱包搬迁」口径不同，不能直接将36.4亿与另一个数字比较来断定谁在说谎。对于金融中介，用户信心既取决于资产负债表，也取决于是否能及时兑付。\cite{ReutersBinance}

\newssub{哪些教训可以迁移到传统软件？}一是密钥与签名权限必须分离，单一工作站被入侵不应足以批准巨额资产转移；二是大额、异常、跨链转账需要独立的速率限制和人工复核；三是事后披露应明确「冻结」「追回」「自行补足」各自数额，避免把资产完整性与攻击溯源混为一谈；四是平台应公布可供审计的储备及负债信息，而不只是展示某一时点的资产截图。技术防线、内部治理与可验证披露是三个层次，缺少其中任何一层都可能让一次入侵变成信任危机。\cite{ReutersBinance,ReutersBinanceCover}
}

\newsitem{Swift望远镜即将再入：一场失败的升轨尝试留下了什么？}{航天工程重点｜NASA10月23日任务结束公告，AP11月1日更新再入预测}{
NASA的尼尔·格雷尔斯Swift天文台2004年11月发射，原定主要任务寿命只有两年，却持续研究伽马射线暴及其X射线、紫外与可见光余辉二十余年。NASA宣布它在10月20日完成最后一次观测，随后在轨道高度约275公里时进入安全模式并结束科学任务。11月1日美联社援引美国太空军预测称，卫星可能在本周晚些时候再入大气层。\textbf{预测的时间和地点存在不确定性}，不能把「本周可能再入」写成已经坠落或必定落在某一国家。\cite{NASASwift,APSwift}

\begin{center}
\resizebox{\linewidth}{!}{\begin{tikzpicture}[>=Latex,x=1cm,y=.75cm,every node/.style={font=\scriptsize}]
\draw[->] (0,0)--(8.2,0) node[right]{时间（示意）};
\draw[->] (0,0)--(0,4) node[above]{轨道高度};
\draw[dashed,gray] (0,1.4)--(7.6,1.4) node[right]{300 km};
\draw[blue!65!black,very thick] (0.5,3.5) .. controls (3.5,3.35) and (4.8,2.8) .. (6.9,.75);
\fill[blue!65!black] (6.9,.75) circle (2pt);
\node[anchor=west] at (4.4,3.4) {太阳活动增强\,$\Rightarrow$\,阻力加大};
\node[anchor=east] at (6.8,.45) {10月20日约275 km};
\end{tikzpicture}}
\par\scriptsize 图3｜Swift轨道衰减的机制示意，\textbf{不是实测时间序列}；275 km来自NASA10月23日公告，曲线与坐标比例均非实测。\end{center}\color{black}

\newssub{为什么卫星会越飞越低？}低地球轨道仍存在极稀薄的大气。太阳活动增强时，高层大气受热膨胀，卫星在同一高度遭遇的阻力增大；阻力消耗轨道机械能，使轨道进一步降低，进而受到更强阻力，形成加速衰减的过程。Swift需要不断改变姿态追踪瞬时天体事件，这种操作有时与减小阻力的最佳姿态冲突。任务团队曾暂停部分科学观测来争取升轨时间。\cite{NASASwift,NASASwiftOps}

\newssub{升轨失败不是简单的「火箭没发射」。}NASA曾委托Katalyst Space的LINK服务飞船尝试靠近并抬升Swift轨道。LINK已经发射，但因自身姿态控制问题，8月的计划不再执行捕获和升轨，转而测试交会与近距离操作能力。对于没有专门为商业机器人对接设计的老卫星，近距离捕获、姿态匹配与结构承载本就存在高风险。NASA称相关试验仍提供了技术经验，不能把它等同于成功挽救Swift。\cite{NASABoost,NASASwift}

\newssub{任务寿命结束，不等于科学价值清零。}Swift最有价值的设计之一，是发现伽马射线暴后快速转向，让多个波段的望远镜追踪余辉。它留下的大量历史观测可供未来重新分析、与其他天文台数据交叉匹配。卫星再入可能结束硬件寿命，却不会自动删除档案。对于未来服务任务，必须比较升轨的成本与失败风险、目标剩余科学寿命，以及服务飞船本身的可靠性。\cite{NASASwift}
}

\newsitem{OPEC+维持12月产量目标：配额不等于能卖出的原油}{能源经济重点｜11月1日七个核心成员会议及路透社报道}{
11月1日，沙特、俄罗斯、伊拉克、科威特、阿尔及利亚、哈萨克斯坦和阿曼七个OPEC+成员决定维持12月石油产量目标不变。这里的「不变」指官方生产配额，而不是所有成员实际日产量已经达到配额。战争扰乱了海湾出口和部分扩产项目，一些生产国即使名义上获准生产更多，仍受到港口、航线、设备和运输保险的约束。\cite{ReutersOPEC}

\newssub{产能、产量、出口量必须分开。}产能是可持续生产的上限，产量是油井实际抽出的原油，出口量则是能离开生产国并进入国际市场的部分。即使产量不变，库存变动和运输中断也可能让出口量发生变化。OPEC+正在评估成员产能，为2027年基准配额作准备；如果高估可用产能，纸面增产可能无法兑现；如果战争结束后运输恢复，供给也可能突然回升并压低价格。\cite{ReutersOPEC}

\newssub{为什么下次会议更关键？}路透社称，全体21个成员计划11月29日讨论2027年框架。七国此前逐步取消一层每日165万桶的自愿减产，另一层约每日200万桶的减产安排仍延续至2026年底。当前讨论的焦点因此不是「今天多抽几桶油」，而是如何在战时供应紧张和战后可能出现过剩之间设定新的基线。11月1日的决定是确定的政策动作，未来油价走向仍是带条件的预测。\cite{ReutersOPEC}
}

\newsitem{法国第三个抗议周末：示威人数下降，社会诉求没有消失}{欧洲社会重点｜10月31日活动、11月1日AP后续报道}{
法国连续第三个周末出现大规模示威。美联社援引法国政府估计，10月17日全国约55万人参与，24日约41万人，31日约32万人；组织者通常给出更高数字。政府口径显示参与人数连续下降，但不能据此直接断言抗议已经结束。运动最初来自高中生对校舍失修、师资短缺和升学机会不均的不满，随后与物价、住房、工资和公共服务诉求汇合。\cite{APFrance,APFranceFollow}

\begin{center}
\begin{tikzpicture}[x=.008cm,y=.65cm,every node/.style={font=\scriptsize}]
\foreach \name/\n/\y in {10月17日/550/2.4,10月24日/410/1.5,10月31日/320/.6}{
\fill[blue!40] (0,\y) rectangle (\n,\y+.38);
\node[anchor=east] at (-5,\y+.2) {\name};
\node[anchor=west] at (\n+6,\y+.2) {\n\,千人};
}
\end{tikzpicture}
\par\scriptsize 图5｜法国政府估计的全国示威人数，单位千人，来源AP 11月1日报道；主办方估计通常更高。不是独立人流测量。\end{center}\color{black}

\newssub{和平集会与街头暴力不是一回事。}报道指出，大多数游行和平进行，但巴黎及其他城市有小群体投掷物品、纵火或袭击警察。11月1日更新的警方统计称，10月31日拘捕779人，低于10月17日的2628人。警方使用催泪瓦斯、水炮及其他驱散手段，也引发比例原则与个人自由的争议。衡量执法是否正当，需要具体事件证据，不能因为少数暴力事件就把数十万示威者全部视为暴徒，也不能因为多数人和平就忽略受害者。\cite{APFranceFollow}

\newssub{德国和葡萄牙的诉求为何相似？}德国多地年轻人要求改善教育、住房和气候政策；法兰克福夜间发生打砸抢，警方报告逾百名警员受伤、54人被拘，主办方与暴力行为划清界限。葡萄牙里斯本也有住房负担抗议，警方清理了一处帐篷营地。各国制度、财政和事件触发点不同，不能将其简化为一场统一组织的「欧洲革命」。但公共服务与生活成本确实构成跨国共同主题。\cite{APFranceFollow,ReutersGermany}

\newssub{接下来观察什么？}法国学校周一复课，工会及学生组织还计划周四进一步行动。应关注教师配置、住房援助和公共财政是否出现具体政策回应，而不是仅以警方拘捕数字衡量政治问题的解决程度。\cite{APFranceFollow}
}

\newsitem{中秘自贸协定升级议定书生效：新增「供应链」章节意味着什么？}{中国对外经济政策重点｜11月1日商务部消息，新华社经中国政府网发布}{
11月1日，中国与秘鲁自由贸易协定升级议定书正式生效。原协定2009年签署、2010年实施；双方2018年启动升级谈判，2024年11月签署升级议定书。此次不是两国第一次建立自贸关系，而是对既有框架的修订：改进原产地规则、海关程序、卫生与植物检疫、知识产权、服务贸易、商务人员临时入境和投资等七个章节，新增标准与合格评定、竞争政策、电子商务、全球供应链及环境与贸易五章。\cite{ChinaPeru}

\newssub{为什么「原产地规则」很重要？}一件商品从多国采购零件、在秘鲁或中国组装，能否享受协定优惠关税，取决于规则如何界定它的经济国籍。规则太复杂，企业即使理论上享有低关税，也可能因证明成本高而放弃优惠；海关程序和标准互认同样会影响交货时间。升级议定书新增供应链章节，鼓励多式联运与物流通道建设，但政策条文生效不等于铁路、港口或仓库已经建成。\cite{ChinaPeru}

\newssub{数字与尚未完成的谈判。}中国官方称双边贸易额从2009年的64亿美元增加到2025年的504亿美元；这是两个年度的名义贸易额比较，不能自动归因于自贸协定一项政策。服务贸易和投资负面清单谈判仍属于未来事项，双方表示将在议定书生效后两年内启动。观察升级效果，应关注实际关税优惠使用率、通关时间和中小企业参与，而非只看协议新增了多少章节。\cite{ChinaPeru}
}

\newsitem{UTe$_2$的两种超导态会「争地盘」：降温反而可能让一种超导消失}{凝聚态物理重点｜《自然·物理》10月30日论文；康奈尔大学研究团队发布实验解读}{
通常说起超导，人们想到的是温度降得足够低，电阻便突然消失。但不同超导态之间可能存在竞争：一种电子配对状态稳定时，另一种反而会被排挤。康奈尔大学Sahas Kamat、Brad Ramshaw等研究者10月30日在《自然·物理》发表论文，报告碲化铀UTe$_2$在低温与高压条件下存在多组分超导行为，并通过热力学测量重新绘制了压力、温度和磁场相关的相图。\cite{UTe2Paper,UTe2Cornell}

\newssub{相图是什么？}它是告诉研究者在什么外界条件下，物质处于什么状态的地图。水的冰、液体和水蒸气可以用温度—压强相图描述；超导体则可能同时受到温度、压力和磁场影响。研究团队利用脉冲超声等手段观察材料的弹性响应，发现先前某些相图解释与热力学约束不一致，于是重新测量，并在极低温及高压下发现不同超导状态相交的区域。\cite{UTe2Cornell}

\newssub{为什么越冷反而会失去一种超导？}这并不是材料整体突然失去所有超导性，而是两种不同对称性的超导序参量发生竞争。当其中一种状态在更低温度下变得更有利，它可能压制原本存在的另一种状态。团队报告在UTe$_2$中观察到这种反直觉的竞争现象。它提示我们：相变不能只按「温度降低，秩序增加」的一条直线来想象；多个序参量之间的耦合可能使相图出现交叉、共存与排斥。\cite{UTe2Paper,UTe2Cornell}

\newssub{离拓扑量子计算还有多远？}某些拓扑超导态理论上可为抗噪声量子信息提供特殊资源，但\textbf{存在多组分超导并不等于已经证明拓扑超导，更不等于做出了可用的拓扑量子比特}。研究者本人也强调拓扑性质尚待直接确认。与网络上常见的「超导量子计算重大突破」标题相比，实验真正扎实的增量，是在压力下厘清了两种超导态的相互作用。论文已发表于同行评议期刊；本刊未独立复现实验。\cite{UTe2Paper,UTe2Cornell}
}

\newsitem{PHP2Shell的后续：修补程序已经发布，真实入侵却可能早已开始}{软件安全重点｜10月29日Unit 42更新、10月30—31日持续报道；原始厂商公告}{
PHP在10月21日公开修复两项严重远程代码执行漏洞CVE-2026-109618和CVE-2026-109619，统称PHP2Shell。10月29日，Palo Alto Networks旗下Unit 42更新安全通告，确认暴露在互联网的PAN-OS管理界面正遭利用，并表示已独立验证Dream Research所报告攻击者使用的工具及结果。Dream声称，攻击者将PHP漏洞与一项Linux提权缺陷串联，在多个国家取得24台防火墙及一台内部管理服务器的最高权限。\textbf{25台服务器是Dream报告的系统数量，不等于25家机构，也不等于所有受害者名单已经公开确认。}\cite{Unit42PHP,DreamSummary}

\newssub{为什么打了补丁还不够？}漏洞修复能阻止同一攻击链继续进入系统，却不保证已经入侵的设备自动恢复可信状态。若攻击者此前已植入持久化组件、窃取管理凭据或建立远程连接，管理员需要检查历史日志、异常账号、证书、设备配置与内部横向移动迹象。尤其是防火墙的管理接口，本不应直接向全互联网开放；Unit 42建议限制到可信内部地址并安装对应PAN-OS修复版本。\cite{Unit42PHP}

\newssub{「AI参与攻击」仍需要证据。}部分报道提到攻击者工具环境出现代码助手及大模型配置，但发现工具痕迹不能证明关键漏洞利用由AI自动完成，更不能据此确认攻击者国籍或政府背景。Dream是一家提供商业安全产品的公司，其事件报告具有技术参考价值，也有商业利益背景；Unit 42确认的是攻击工具与相关结果，并未独立公开全部受害机构清单。当前最可靠的运维结论是：\textbf{更新、隔离管理面、审计已失陷系统}，而不是根据未经核实的归因传闻寻找替罪者。\cite{Unit42PHP,DreamSummary}
}

\newsitem{数学拾趣｜锦标赛的两个部分都「很复杂」，为何仍找不到指定的两组三角形？}{组合数学｜Chen、van der Holst、Wang；arXiv:2610.37920，10月27日提交、30日新列表；未同行评议}{
数学里的\textbf{锦标赛}不是体育赛事，而是把完全图的每条边都指定一个方向后得到的有向图。任意两个顶点之间恰有一条箭头；三个顶点若构成$a\to b\to c\to a$，就形成\textbf{有向循环三角形}。如果一个锦标赛能拆成少数几个没有有向循环的部分，便说它的锦标赛色数较小；需要很多部分才能做到，则色数很大。\cite{TournamentPaper}

\begin{center}
\begin{tikzpicture}[>=Stealth,every node/.style={font=\scriptsize}]
\foreach \a/\x/\y in {a/0/1.5,b/-1.2/-.3,c/1.2/-.3}{\node[draw,circle,fill=blue!10,minimum size=6mm] (\a) at (\x,\y) {\a};}
\draw[->,thick] (a)--(b); \draw[->,thick] (b)--(c); \draw[->,thick] (c)--(a);
\node[align=left] at (4.1,.6) {有向三角形：$a\to b\to c\to a$\\每对顶点之间恰有一条定向边\\这是锦标赛中的一个循环};
\end{tikzpicture}
\par\scriptsize 图4｜依据图论定义自行绘制的三顶点例子，不是论文反例的实际大型构造。\end{center}\color{black}

\newssub{从定义到困难。}Nguyen、Scott和Seymour曾猜想：当一个锦标赛分成两个互不相交、各自色数都足够大的部分时，应该能在每部分各找到一个循环三角形，并且这六个顶点之间跨部分的九条边方向一致。直觉上，两个部分都极其复杂，似乎应当产生足够丰富的局部结构；但「分别存在复杂结构」并不自动保证「两者之间满足统一方向」。\cite{TournamentPaper}

\newssub{作者的新方法。}Guantao Chen、Hein van der Holst和Shengze Wang在预印本中声称构造出反例：对于任意给定正整数$d$，都存在一个锦标赛及划分$(X,Y)$，使两个部分的色数都至少为$d$，却找不到满足跨部分统一定向的循环三角形对。构造结合大围长、大色数的图以及二进制函数版本的Kneser图；色数下界给出两种证明路线，分别涉及Borsuk--Ulam定理与Tucker引理。\cite{TournamentPaper}

\newssub{验证状态。}这是作者10月27日提交的arXiv预印本，出现在10月30日的最新分类列表。它的「推翻猜想」目前是作者提供的证明主张；本刊没有独立重建反例、逐项验证所有参数，也没有证据表明已完成期刊同行评议。此题有趣之处在于提醒我们：大色数保证某种整体复杂性，却不必强迫两个复杂子系统之间出现预期的局部配置。\cite{TournamentPaper}
}

\newsitem{数学简讯｜一个六十年代的定理可能需要撤回}{集合论与序理论｜Benhamou、Jakob；arXiv:2610.38947，10月28日提交；未同行评议}{
Tom Benhamou和Hannes Jakob发表23页预印本《Stationary Limits of the Point Spectrum》，声称构造出一个有向集，其Tukey谱在平稳极限下不封闭，从而给出Isbell在20世纪60年代提出的一项定理的反例。这里的「谱」不是物理光谱，而是用来刻画有向集在Tukey约化意义下的结构类型。论文进一步提出介于平稳反射与弱紧性之间的组合原则。作者明确欢迎评论；本刊尚未检查其集合论细节，也未确认独立复核结果，故不将其作为已获学界确认的纠错结论。\cite{TukeyPaper}
}

\newsitem{航天简讯｜模拟火星任务四名成员出舱，但不能代表真实火星航行}{NASA与人因工程｜11月1日路透社现场报道；与第019期预告相比新增实际结束确认}{
路透社11月1日报道，NASA第二次CHAPEA模拟火星任务的四名成员在休斯敦约378天的隔离生活后走出栖居舱。这补上了前一期仅能确认NASA预定出舱时间、尚待核实是否完成的证据缺口。模拟任务可观察隔离、工作安排、通信延迟与资源限制下的团队表现，却无法完整再现火星重力、辐射和不可快速撤离的风险。实际出舱已得到现场报道支持；科研数据的系统分析、公开发表与独立评议仍是后续事项。\cite{ReutersCHAPEA}
}

\par\medskip{\large\bfseries\color{newsblue}来源、图像与核查说明}\par
{\scriptsize 本期主动检查国际时政外交、战争、国内政策、经济、社会议题、AI与网络安全、Rust/C++、Linux、物理、化学、航天、文化及arXiv数论、代数、几何、组合、逻辑、概率、Lean/mathlib。未有可靠且足够新的材料时不强行刊登。五张TikZ矢量图均由编辑部依据已核查来源或数学定义绘制，并以实际源文件归档。NASA Swift原始照片及AP现场照片可从来源页核查，但本次未能取得可合法归档的图像文件；因此未使用网页缩略图，也未使用AI合成新闻现场照片。新闻发生日期与报道发表日期已分开；对来源立场、未决责任和预印本核验状态均作保留。}\par
\begingroup\fontsize{6}{7.2}\selectfont
\renewcommand{\refname}{参考文献}
\bibliographystyle{plainurl}\bibliography{refs}
\endgroup
\end{documents}
\end{document}